\SourceLeaf{59}
\begin{center}
{\large NEW PROOFS AND EXTENSIONS\par}
\smallskip OF THE\par
\smallskip{\large FUNDAMENTAL THEOREM\par}
\medskip IN THE THEORY OF QUADRATIC RESIDUES.\par
\end{center}
\BookRule{20mm}
The fundamental theorem about quadratic residues, which ranks among the most beautiful truths of higher arithmetic, was readily discovered by induction, but proved with far greater difficulty. It often happens in this field that the proofs of the simplest truths, which seem to offer themselves to the investigator through induction of their own accord, lie deeply hidden and can finally be brought to light only after many unsuccessful attempts, perhaps by a path very different from the one along which they were sought. Then it often happens that, once one path has been found, \textit{several} others gradually open up, leading to the same goal: some more briefly and directly, others as if by a roundabout way, starting from very different principles between which and the question at hand one would scarcely have suspected any connection. This remarkable connection among more hidden truths not only gives these investigations a particular charm; it also deserves careful study and unravelling, since it often yields new resources or advances for the science itself.

Although the arithmetical theorem considered here might therefore seem, through the earlier efforts that supplied four entirely different proofs\begingroup\renewcommand{\thefootnote}{*)}\footnote{Two are given in the fourth and fifth sections of the \textit{Disquisitiones Arithmeticae}; the third in a separate paper (\textit{Commentt. Soc. Gotting. Vol. XVI}), and the fourth in the paper \textit{Summatio quarundam serierum singularium} (\textit{Commentt. Recentiores, Vol. I}).}\endgroup, to have been fully\SourceLeaf{60} dealt with, I nevertheless return to the same subject and add two further proofs, which will certainly shed new light on it. The first is related to the third in a certain respect, since it starts from the same lemma; it then follows a different path, however, so that it can rightly be regarded as a new proof, which in elegance will seem, if not superior to that third proof, at least by no means inferior. The sixth proof rests on an entirely different and subtler principle, and gives a new example of the remarkable connection between arithmetical truths that at first sight seem very far apart. These two proofs are accompanied by a very simple new algorithm for deciding whether a given integer is a quadratic residue or a nonresidue of a given prime number.

\newpage
There was another reason that led me to publish now, in particular, the new proofs promised nine years ago. When I began in 1805 to investigate the theory of cubic and biquadratic residues, a far more difficult subject, I experienced much the same course of events as before in the theory of quadratic residues. The theorems that completely settle these questions, and display a remarkable analogy with the theorems about quadratic residues, were immediately discovered by induction as soon as they were sought along a suitable path; but all attempts to obtain proofs complete in every respect remained unsuccessful for a long time. This was precisely what prompted me to strive so eagerly to add more and more proofs to those already known for quadratic residues, sustained by the hope that, among many different methods, one or another might help to illuminate the related subject. This hope was by no means in vain, and tireless work was finally followed by good results. Soon it will be possible to bring the fruits of my nights of work before the public; but before I undertake this arduous task, I have resolved to return once more to the theory of quadratic residues, complete everything that still remains to be done in it, and thus, as it were, take leave of this part of higher arithmetic.
